\chapter*{TÓM TẮT NỘI DUNG ĐỒ ÁN}

UniTrack là nền tảng Web hỗ trợ giảng viên giám sát quá trình thực hiện dự án của sinh viên. Hệ thống tập trung vào quy trình dự án trước tiên: giảng viên tạo và quản lý dự án, phân công nhiệm vụ chính thức, theo dõi tiến độ, xem minh chứng, phản hồi và duyệt kết quả nộp bài của sinh viên. Sinh viên sử dụng hệ thống để theo dõi công việc được giao, gửi báo cáo tiến độ, bổ sung tài nguyên hoặc tệp minh chứng và nhận phản hồi từ giảng viên.

Đồ án này trình bày quá trình khảo sát yêu cầu, thiết kế kiến trúc, xây dựng và kiểm thử hệ thống UniTrack. Ứng dụng được phát triển theo mô hình Web hiện đại với frontend React/TypeScript, backend Go, cơ sở dữ liệu PostgreSQL, API REST và cơ chế phiên đăng nhập bằng cookie. Hệ thống cũng chú trọng tới phân quyền theo vai trò, kiểm soát truy cập dự án, ràng buộc toàn vẹn dữ liệu và các luồng xử lý an toàn cho tài khoản, dự án, nhiệm vụ, nộp bài, tài nguyên và tệp minh chứng.

Kết quả của đồ án là một sản phẩm có thể chạy cục bộ, có cấu trúc mã nguồn rõ ràng, tài liệu theo từng tính năng và bộ kiểm thử backend/frontend cho các luồng quan trọng. Báo cáo này tổng hợp mục tiêu, phạm vi, công nghệ sử dụng, thiết kế hệ thống, kết quả triển khai, kiểm thử, các đóng góp nổi bật và hướng phát triển tiếp theo của UniTrack.

\vspace{1cm}
\begin{flushright}
Sinh viên thực hiện

(Ký và ghi rõ họ tên)
\end{flushright}
